% Generated draft: requires mathematical and source review.
\documentclass[11pt]{article}
\usepackage[margin=25mm]{geometry}
\usepackage{fontspec}
\setmainfont{Latin Modern Roman}
\usepackage{amsmath,amssymb,mathtools,unicode-math}
\setmathfont{Latin Modern Math}
\usepackage{enumitem,needspace,xurl,hyperref}
\setlength{\parindent}{0pt}
\setlength{\parskip}{5pt}
\setlength{\emergencystretch}{3em}
\newcommand{\N}{\mathbb{N}}
\newcommand{\Z}{\mathbb{Z}}
\newcommand{\Q}{\mathbb{Q}}
\newcommand{\R}{\mathbb{R}}
\newcommand{\C}{\mathbb{C}}
\newcommand{\F}{\mathbb{F}}
\newcommand{\A}{\mathbb{A}}
\newcommand{\PP}{\mathbb{P}}
\newcommand{\E}{\mathbb{E}}
\newcommand{\eps}{\varepsilon}
\newcommand{\dd}{\,\mathrm{d}}
\begin{document}
\section*{UnsolvedMath ID: 10600017}
\textbf{Source problem code:} catalogue entry
\subsection{Definitions and mathematical statement}
For a virtual link \(K\), let \(\mathcal D(K)\) be its diagrams modulo virtual-link equivalence. If \(c(D)\) and \(v(D)\) count classical and virtual crossings of \(D\), respectively, define
\[
 C(K)=\min_{D\in\mathcal D(K)}c(D),\quad
 V(K)=\min_{D\in\mathcal D(K)}v(D),\quad
 T(K)=\min_{D\in\mathcal D(K)}(c(D)+v(D)).
\]
These minima are taken independently; a diagram minimizing one need not minimize the others.

For a virtual knot, the source describes classical unknotting number by minimizing the number of classical crossing switches needed to reach the unknot, over all diagrams. If such switches cannot unknot it, an infinite/undefined-value convention is needed. The virtual unknotting number instead minimizes the number of classical crossings changed directly into virtual crossings, again allowing any virtual diagram before the changes. This is the operation explicitly described in the source as ``direct flattening''; it must not be replaced by a different virtualization move or by merely forgetting over/under information at a flat crossing.

Problem-list item 17 retains the following requests:
\begin{enumerate}
\item Investigate the virtual unknotting number, beyond the source's statement that very little is known.
\item Relate \(V(K)\) to the least genus of a surface representation of \(K\). In the usual surface model the latter minimizes genus over representations in thickened oriented surfaces up to the intended stabilization equivalence, but the source gives no exact conventions.
\item Decide the printed equality \(T=V+L\). The source has defined \(C,V,T\), while \(L\) denotes a link and has not been defined as a numerical invariant. Thus this equality is ill-typed as extracted. If \(L\) was intended to mean \(C\), the candidate becomes \(T(K)=V(K)+C(K)\), asking for simultaneous attainment of the two crossing minima. This is an explicit conditional reading, not a silent correction.
\item Find an algorithm for \(V\) for some class of virtual knots. The class and effective input model are unspecified.
\item Obtain nontrivial upper and lower bounds from virtual-knot polynomials for the printed list \(T,V,L\). The third numerical quantity \(L\), the polynomial families and the meaning of nontrivial bounds are not supplied.
\end{enumerate}
The source reports partial work for \(T\) in two link classes, one called quasialternating and one unnamed, and mentions classical alternating diagrams; their citation slots are empty. The supplied modern summary reports arrow/writhe-polynomial lower bounds and special virtual-torus-knot unknotting computations, while saying the printed equality and algorithmic \(V\) requests remain unresolved. Its scope and references require review.
\subsection{Short English statement}
Study virtual unknotting, compare minimal virtual crossings with surface genus, develop algorithms and polynomial bounds, and clarify the printed crossing-number equality. The extract defines \(C,V,T\) but later uses an undefined numerical \(L\); this is recorded instead of silently replacing it by \(C\). The supplied progress assessment has incomplete citations.
\subsection{Sources}
\begin{itemize}
\item Local UnsolvedMath registry, permanent ID \(10600017\), exact source code \texttt{catalogue entry}, ``Virtual-knot problem 17 --- Crossing number problems.'' No individual proposer or proposed year is supplied.
\item \emph{Unsolved Problems in Virtual Knot Theory and Combinatorial Knot Theory}, revised survey, dedicated problem-list item 17: \url{https://arxiv.org/abs/1409.2823}. Background: 60 top-level items, compound items retained intact; source-TeX extraction, accessed 2026-07-29. The in-statement references for partial algorithms, alternating diagrams and polynomials are empty.
\item Supplied incomplete references, identifiers retained literally: ``Lower bounds on virtual crossing number and minimal surface genus'' (2009), arXiv:0904.xxxx; ``Virtual Crossing Number and the Arrow Polynomial'' (2008), arXiv:0810.xxxx; ``On Virtual Crossing Number Estimates For Virtual Links'' (2008), arXiv:0811.xxxx.
\item Further supplied incomplete references: ``On Virtual Crossing Numbers for Virtual Knots'' (2011), arXiv:1107.xxxx; ``A note on the writhe polynomial and the virtual crossing number'' (2018), arXiv:1805.xxxx; ``Virtual unknotting numbers of certain virtual torus knots'' (2017), arXiv:1701.xxxx. These xxxx identifiers are placeholders, not verified retrievable citations.
\item Registry status partially\_solved and classification PARTIAL-PROGRESS report polynomial bounds and computations for special classes, but no resolution found for \(T=V+L\) or general \(V\) algorithms. Source evidence NEEDS\_REVIEW presents an open problem and calls for a current primary-literature audit.
\item Dataset attribution: Ulam AI / UnsolvedMath Contributors, CC BY 4.0. Source rights NEEDS\_REVIEW: publicly accessible author/publisher copy, redistribution terms requiring review.
\end{itemize}
\end{document}
